% ============================================================
\chapter{S\'eparation de Variables et S\'eries de Fourier}
\label{ch:fourier}
% ============================================================

En 1807, Joseph Fourier pr\'esente \`a l'Acad\'emie des Sciences de Paris un m\'emoire audacieux~: toute fonction p\'eriodique, aussi irr\'eguli\`ere soit-elle, peut s'\'ecrire comme une somme de sinus et de cosinus. L'affirmation est accueillie avec scepticisme --- Lagrange la juge impossible --- mais elle se r\'ev\'elera \^etre l'une des id\'ees les plus f\'econdes de toute l'histoire des math\'ematiques. La m\'ethode de Fourier, n\'ee de l'\'etude de la propagation de la chaleur, offre une strat\'egie remarquable pour r\'esoudre les EDP lin\'eaires~: d\'ecomposer le probl\`eme en modes propres, r\'esoudre chaque mode s\'epar\'ement, puis recombiner. Ce principe de \emph{s\'eparation des variables} transforme une EDP en une famille d'EDO ordinaires --- un gain consid\'erable.

\section{Principe de la s\'eparation des variables}

La m\'ethode de s\'eparation des variables (ou m\'ethode de Fourier) est l'une
des techniques les plus puissantes pour r\'esoudre les EDP lin\'eaires sur des
domaines simples avec des conditions aux limites homog\`enes.

\begin{intuition}
L'id\'ee est de chercher la solution sous la forme d'un produit de fonctions,
chacune ne d\'ependant que d'une seule variable. L'EDP se d\'ecompose alors
en plusieurs EDO, beaucoup plus faciles \`a r\'esoudre.
\end{intuition}

\subsection{M\'ethode g\'en\'erale}

Consid\'erons une EDP de la forme $L[u] = 0$ pos\'ee sur un domaine
$\Omega = (0,L) \times (0,T)$ avec conditions aux limites homog\`enes.

\textbf{\'Etape 1~:} Chercher des solutions de la forme $u(x,t) = X(x)\, T(t)$.

\textbf{\'Etape 2~:} Substituer dans l'EDP et s\'eparer les variables
pour obtenir deux EDO coupl\'ees par une constante de s\'eparation $\lambda$.

\textbf{\'Etape 3~:} R\'esoudre le probl\`eme de Sturm-Liouville pour $X$
(avec conditions aux limites), ce qui d\'etermine les valeurs propres $\lambda_n$.

\textbf{\'Etape 4~:} R\'esoudre l'EDO pour $T$ avec chaque $\lambda_n$.

\textbf{\'Etape 5~:} Former la solution g\'en\'erale par superposition~:
$u = \sum_n c_n X_n(x)\, T_n(t)$.

\textbf{\'Etape 6~:} D\'eterminer les coefficients $c_n$ par la condition initiale
en utilisant les s\'eries de Fourier.

\begin{exemple}[\'Equation de la chaleur sur un segment]
Consid\'erons~:
\[
    \begin{cases}
        u_t = k\, u_{xx}, & 0 < x < L,\; t > 0,\\
        u(0,t) = u(L,t) = 0, & t > 0,\\
        u(x,0) = u_0(x), & 0 < x < L.
    \end{cases}
\]
On pose $u(x,t) = X(x)\, T(t)$. L'EDP donne $X\, T' = k\, X''\, T$, soit~:
\[
    \frac{T'}{kT} = \frac{X''}{X} = -\lambda \quad (\text{constante}).
\]
Ceci donne deux EDO~:
\[
    X'' + \lambda X = 0, \quad X(0) = X(L) = 0,
\]
\[
    T' + k\lambda T = 0.
\]
Le probl\`eme pour $X$ est un probl\`eme de Sturm-Liouville dont les
solutions non triviales sont~:
\[
    \lambda_n = \left(\frac{n\pi}{L}\right)^2, \quad
    X_n(x) = \sin\!\left(\frac{n\pi x}{L}\right), \quad n \geq 1.
\]
L'EDO pour $T$ donne $T_n(t) = e^{-k\lambda_n t}$. La solution g\'en\'erale est~:
\[
    u(x,t) = \sum_{n=1}^{\infty} b_n \sin\!\left(\frac{n\pi x}{L}\right)
    e^{-k(n\pi/L)^2 t},
\]
o\`u les $b_n$ sont les coefficients de Fourier en sinus de $u_0$.
\end{exemple}

\section{S\'eries de Fourier}

\subsection{D\'efinitions et coefficients}

\begin{definition}[S\'erie de Fourier]
Soit $f : [-\pi, \pi] \to \R$ une fonction int\'egrable. La \textbf{s\'erie
de Fourier} de $f$ est~:
\begin{equation}\label{eq:serie_fourier}
    f(x) \sim \frac{a_0}{2} + \sum_{n=1}^{\infty}
    \left[a_n \cos(nx) + b_n \sin(nx)\right],
\end{equation}
o\`u les \textbf{coefficients de Fourier} sont~:
\begin{equation}\label{eq:coefficients_fourier}
    a_n = \frac{1}{\pi}\int_{-\pi}^{\pi} f(x)\cos(nx)\, dx, \quad
    b_n = \frac{1}{\pi}\int_{-\pi}^{\pi} f(x)\sin(nx)\, dx.
\end{equation}
\end{definition}

\begin{definition}[S\'erie de Fourier complexe]
De mani\`ere \'equivalente, on peut \'ecrire~:
\[
    f(x) \sim \sum_{n=-\infty}^{+\infty} c_n\, e^{inx}, \quad
    c_n = \frac{1}{2\pi}\int_{-\pi}^{\pi} f(x)\, e^{-inx}\, dx.
\]
\end{definition}

\begin{formules}[title=\textbf{Coefficients de Fourier sur $[0,L]$}]
Pour une fonction $f$ sur $[0,L]$~:
\begin{itemize}
    \item Coefficients en sinus~: $b_n = \frac{2}{L}\int_0^L f(x)\sin\!\left(\frac{n\pi x}{L}\right) dx$.
    \item Coefficients en cosinus~: $a_n = \frac{2}{L}\int_0^L f(x)\cos\!\left(\frac{n\pi x}{L}\right) dx$.
\end{itemize}
\end{formules}

\subsection{Orthogonalit\'e}

\begin{proposition}[Orthogonalit\'e des fonctions trigonom\'etriques]
Le syst\`eme $\{1, \cos(nx), \sin(nx)\}_{n \geq 1}$ est orthogonal dans
$L^2(-\pi, \pi)$~:
\begin{align*}
    \inner{\cos(mx)}{\cos(nx)} &= \int_{-\pi}^{\pi} \cos(mx)\cos(nx)\, dx
    = \pi\, \delta_{mn}, \\
    \inner{\sin(mx)}{\sin(nx)} &= \int_{-\pi}^{\pi} \sin(mx)\sin(nx)\, dx
    = \pi\, \delta_{mn}, \\
    \inner{\cos(mx)}{\sin(nx)} &= 0 \quad \forall m, n.
\end{align*}
\end{proposition}

\subsection{Convergence des s\'eries de Fourier}

\begin{theoreme}[Convergence ponctuelle --- Dirichlet]
Si $f$ est $2\pi$-p\'eriodique, continue par morceaux et de classe $C^1$
par morceaux, alors la s\'erie de Fourier converge en tout point~:
\[
    \frac{a_0}{2} + \sum_{n=1}^{\infty} (a_n \cos(nx) + b_n \sin(nx))
    = \frac{f(x^+) + f(x^-)}{2}.
\]
En particulier, aux points de continuit\'e, la s\'erie converge vers $f(x)$.
\end{theoreme}

\begin{theoreme}[Convergence uniforme]
Si $f$ est $2\pi$-p\'eriodique, continue et de classe $C^1$ par morceaux,
alors la s\'erie de Fourier converge uniform\'ement vers $f$.
\end{theoreme}

\begin{theoreme}[Convergence dans $L^2$ --- Parseval]\label{thm:parseval}
Si $f \in L^2(-\pi, \pi)$, la s\'erie de Fourier converge vers $f$ dans
$L^2$ et l'\'egalit\'e de \textbf{Parseval} est v\'erifi\'ee~:
\begin{equation}\label{eq:parseval}
    \frac{1}{\pi}\int_{-\pi}^{\pi} \abs{f(x)}^2\, dx
    = \frac{a_0^2}{2} + \sum_{n=1}^{\infty} (a_n^2 + b_n^2).
\end{equation}
\end{theoreme}

\begin{remarque}
L'\'egalit\'e de Parseval exprime que le syst\`eme trigonom\'etrique forme
une \textbf{base hilbertienne} de $L^2(-\pi,\pi)$. C'est un cas particulier
du th\'eor\`eme de compl\'etude.
\end{remarque}

\begin{exemple}
Calculons la s\'erie de Fourier de $f(x) = x$ sur $[-\pi, \pi]$.
Par imparit\'e, $a_n = 0$ pour tout $n \geq 0$. Pour les $b_n$~:
\[
    b_n = \frac{1}{\pi}\int_{-\pi}^{\pi} x \sin(nx)\, dx
    = \frac{2(-1)^{n+1}}{n}.
\]
Donc~:
\[
    x = 2\sum_{n=1}^{\infty} \frac{(-1)^{n+1}}{n}\sin(nx),
    \quad x \in (-\pi, \pi).
\]
L'\'egalit\'e de Parseval donne $\sum_{n=1}^{\infty} \frac{1}{n^2} = \frac{\pi^2}{6}$.
\end{exemple}

\section{Probl\`emes de Sturm-Liouville}

\begin{definition}[Probl\`eme de Sturm-Liouville]
Un \textbf{probl\`eme de Sturm-Liouville r\'egulier} est un probl\`eme aux
valeurs propres de la forme~:
\begin{equation}\label{eq:sturm_liouville}
    -(p(x)\, y')' + q(x)\, y = \lambda\, w(x)\, y, \quad x \in (a,b),
\end{equation}
avec des conditions aux limites s\'epar\'ees, o\`u $p > 0$, $w > 0$ et $q$
sont continues sur $[a,b]$.
\end{definition}

\begin{theoreme}[Propri\'et\'es de Sturm-Liouville]
Sous les hypoth\`eses ci-dessus~:
\begin{enumerate}
    \item Il existe une suite infinie de valeurs propres r\'eelles
          $\lambda_1 < \lambda_2 < \cdots \to +\infty$.
    \item Les fonctions propres $\phi_n$ correspondantes sont orthogonales
          dans $L^2_w(a,b)$~:
          $\int_a^b \phi_m(x)\, \phi_n(x)\, w(x)\, dx = 0$ pour $m \neq n$.
    \item La fonction propre $\phi_n$ admet exactement $n-1$ z\'eros dans
          $(a,b)$.
    \item Le syst\`eme $\{\phi_n\}$ forme une base hilbertienne de $L^2_w(a,b)$.
\end{enumerate}
\end{theoreme}

\begin{exemple}[Probl\`emes types]
\leavevmode
\begin{enumerate}
    \item $X'' + \lambda X = 0$, $X(0) = X(L) = 0$ (Dirichlet)~:\\
          $\lambda_n = (n\pi/L)^2$, $X_n = \sin(n\pi x/L)$.
    \item $X'' + \lambda X = 0$, $X'(0) = X'(L) = 0$ (Neumann)~:\\
          $\lambda_n = (n\pi/L)^2$ pour $n \geq 0$, $X_0 = 1$,
          $X_n = \cos(n\pi x/L)$.
    \item $X'' + \lambda X = 0$, $X(0) = 0$, $X'(L) = 0$ (mixte)~:\\
          $\lambda_n = ((2n-1)\pi/(2L))^2$, $X_n = \sin((2n-1)\pi x/(2L))$.
\end{enumerate}
\end{exemple}

\section{Applications aux EDP}

\subsection{\'Equation de la chaleur avec conditions de Neumann}

\begin{exemple}
R\'esoudre~:
\[
    \begin{cases}
        u_t = k\, u_{xx}, & 0 < x < L,\; t > 0,\\
        u_x(0,t) = u_x(L,t) = 0, & t > 0,\\
        u(x,0) = u_0(x).
    \end{cases}
\]
La s\'eparation donne $X'' + \lambda X = 0$ avec $X'(0) = X'(L) = 0$
(Neumann). Les valeurs propres sont $\lambda_n = (n\pi/L)^2$,
$n \geq 0$, et~:
\[
    u(x,t) = \frac{a_0}{2} + \sum_{n=1}^{\infty} a_n
    \cos\!\left(\frac{n\pi x}{L}\right) e^{-k(n\pi/L)^2 t},
\]
o\`u $a_n = \frac{2}{L}\int_0^L u_0(x)\cos(n\pi x/L)\, dx$.

Quand $t \to \infty$, $u(x,t) \to a_0/2 = \frac{1}{L}\int_0^L u_0(x)\, dx$,
la temp\'erature moyenne initiale.
\end{exemple}

\subsection{\'Equation des ondes}

\begin{exemple}
R\'esoudre~:
\[
    \begin{cases}
        u_{tt} = c^2 u_{xx}, & 0 < x < L,\; t > 0,\\
        u(0,t) = u(L,t) = 0,\\
        u(x,0) = u_0(x), \quad u_t(x,0) = v_0(x).
    \end{cases}
\]
On cherche $u(x,t) = X(x)\, T(t)$, ce qui donne~:
\[
    X'' + \lambda X = 0, \; X(0)=X(L)=0 \quad\text{et}\quad
    T'' + c^2\lambda T = 0.
\]
Les solutions sont $X_n = \sin(n\pi x/L)$,
$\omega_n = cn\pi/L$, et~:
\[
    u(x,t) = \sum_{n=1}^{\infty} \sin\!\left(\frac{n\pi x}{L}\right)
    \left[A_n \cos(\omega_n t) + B_n \sin(\omega_n t)\right],
\]
avec~:
\[
    A_n = \frac{2}{L}\int_0^L u_0(x)\sin\!\left(\frac{n\pi x}{L}\right) dx,
    \quad
    B_n = \frac{2}{L\omega_n}\int_0^L v_0(x)\sin\!\left(\frac{n\pi x}{L}\right) dx.
\]
\end{exemple}

\subsection{\'Equation de Laplace sur un rectangle}

\begin{exemple}
R\'esoudre~:
\[
    \begin{cases}
        \Delta u = 0, & 0 < x < a,\; 0 < y < b,\\
        u(0,y) = u(a,y) = 0,\\
        u(x,0) = 0, \quad u(x,b) = f(x).
    \end{cases}
\]
En posant $u = X(x)\, Y(y)$~:
\[
    \frac{X''}{X} = -\frac{Y''}{Y} = -\lambda.
\]
Avec $X(0) = X(a) = 0$~: $\lambda_n = (n\pi/a)^2$,
$X_n = \sin(n\pi x/a)$. L'\'equation pour $Y$ donne
$Y'' - \lambda_n Y = 0$ avec $Y(0) = 0$~:
\[
    Y_n(y) = \sinh\!\left(\frac{n\pi y}{a}\right).
\]
La solution est~:
\[
    u(x,y) = \sum_{n=1}^{\infty} c_n \sin\!\left(\frac{n\pi x}{a}\right)
    \sinh\!\left(\frac{n\pi y}{a}\right),
\]
avec $c_n = \frac{2}{a\sinh(n\pi b/a)}\int_0^a f(x)\sin(n\pi x/a)\, dx$.
\end{exemple}

\section{S\'eparation en coordonn\'ees polaires}

\begin{exemple}[Laplace dans le disque]
R\'esoudre $\Delta u = 0$ dans le disque $r < R$ avec $u(R,\theta) = f(\theta)$.
En polaires, $\Delta u = u_{rr} + \frac{1}{r}u_r + \frac{1}{r^2}u_{\theta\theta} = 0$.
On pose $u = \mathcal{R}(r)\, \Theta(\theta)$~:
\[
    r^2\frac{\mathcal{R}''}{\mathcal{R}} + r\frac{\mathcal{R}'}{\mathcal{R}}
    = -\frac{\Theta''}{\Theta} = \lambda.
\]
La p\'eriodicit\'e en $\theta$ impose $\lambda_n = n^2$, $n \geq 0$.
Les solutions sont~:
\[
    u(r,\theta) = \frac{a_0}{2} + \sum_{n=1}^{\infty}
    \left(\frac{r}{R}\right)^n \bigl[a_n \cos(n\theta) + b_n \sin(n\theta)\bigr],
\]
o\`u $a_n$, $b_n$ sont les coefficients de Fourier de $f$.
Ceci donne la \textbf{formule int\'egrale de Poisson}~:
\[
    u(r,\theta) = \frac{1}{2\pi}\int_0^{2\pi}
    \frac{R^2 - r^2}{R^2 - 2Rr\cos(\theta-\phi) + r^2}\, f(\phi)\, d\phi.
\]
\end{exemple}

\section{Exercices}

\begin{exercice}
Calculer la s\'erie de Fourier de $f(x) = \abs{x}$ sur $[-\pi,\pi]$.
En d\'eduire $\sum_{k=0}^{\infty} \frac{1}{(2k+1)^2} = \frac{\pi^2}{8}$.
\end{exercice}

\begin{exercice}
R\'esoudre l'\'equation de la chaleur $u_t = u_{xx}$ sur $(0,\pi)$
avec $u(0,t) = u(\pi,t) = 0$ et $u(x,0) = \sin(x) + 3\sin(2x)$.
\end{exercice}

\begin{exercice}
R\'esoudre l'\'equation des ondes $u_{tt} = 4u_{xx}$ sur $(0,1)$
avec $u(0,t) = u(1,t) = 0$, $u(x,0) = x(1-x)$, $u_t(x,0) = 0$.
\end{exercice}

\begin{exercice}
R\'esoudre le probl\`eme de Laplace sur le rectangle $(0,\pi) \times (0,1)$
avec $u(0,y) = u(\pi,y) = 0$, $u(x,0) = 0$, $u(x,1) = \sin(3x)$.
\end{exercice}

\begin{exercice}
Montrer que les coefficients de Fourier de $f$ satisfont
$a_n, b_n \to 0$ quand $n \to \infty$ (lemme de Riemann-Lebesgue).
\end{exercice}

\begin{exercice}
R\'esoudre le probl\`eme de Dirichlet dans le disque $r < 1$
avec $u(1,\theta) = \cos^2\theta$. \'Ecrire la solution sous
forme ferm\'ee.
\end{exercice}

\begin{exercice}[Probl\`eme de Sturm-Liouville singulier]
Trouver les valeurs propres et fonctions propres de~:
\[
    -\frac{d}{dx}\left(x\frac{dy}{dx}\right) = \frac{\lambda}{x}\, y,
    \quad 1 < x < e, \quad y(1) = y(e) = 0.
\]
\emph{Indication~:} poser $x = e^s$.
\end{exercice}
